% ============================================================
% section5_processing.tex - Раздел 5: Цифровая обработка сигналов
% ============================================================
% !TeX program = xelatex
% Компиляция: xelatex --shell-escape main.tex

\section{Цифровая обработка: Окна и спектры}\label{sec:processing}

В предыдущем разделе мы успешно загрузили и структурировали наши временные ТГц-импульсы. Однако фундаментальные физические свойства исследуемой проволочной решетки (такие как частотный коэффициент пропускания по мощности) проявляются в частотной (спектральной) области. 

Для перехода из временной области в частотную в цифровой обработке сигналов (DSP, от англ. Digital Signal Processing --- математические методы анализа и преобразования дискретных сигналов) используется классический инструмент — Быстрое преобразование Фурье (БПФ, или FFT, от англ. Fast Fourier Transform --- быстрый алгоритм спектрального разложения). Наша задача — использовать БПФ как готовый инженерный рецепт для превращения зависимости напряженности поля от времени в зависимость амплитуды от частоты.

\note{\textbf{Математическая теория спектрального анализа (для любознательных)}\\
В этой книге мы не будем углубляться в строгое математическое доказательство рядов и интегралов Фурье, свойства ортогональности гармонических функций и теоремы дискретизации. Мы рассматриваем БПФ исключительно как практический инструмент. Если ты хочешь глубоко изучить теоретические основы спектрального анализа и цифровой фильтрации, мы рекомендуем обратиться к классическим учебникам, например: \textbf{А. Оппенгейм, Р. Шафер. Цифровая обработка сигналов} или \textbf{С. Смит. Цифровая обработка сигналов. Практическое руководство для инженеров и научных работников}.}

Давай разберем пошаговый рецепт подготовки сигнала и расчета спектра, а затем оформим его в виде готового модуля \mycode{spectrum.py}.

\subsection{Формат данных и типы в NumPy: Знакомимся с np.ndarray}

Прежде чем приступать к цифровой обработке, давай вспомним, в каком именно формате возвращает данные наша функция \mycode{load_tds_data}, написанная в Разделе~\ref{sec:loading} (Глава~3). При вызове \mycode{t, E = load_tds_data(filepath)} мы получаем два связанных одномерных числовых вектора:
\begin{enumerate}
    \item \mycode{t} (массив временных отсчетов) — представляет собой сетку времени в пикосекундах ($\text{пс}$), обычно с шагом порядка $0.05$--$0.1$~пс, общей длительностью записи около $100$~пс.
    \item \mycode{E} (массив амплитуды поля) — содержит соответствующие значения напряженности электрического поля терагерцового импульса в условных единицах (у.е.).
\end{enumerate}

Оба этих объекта имеют тип \mycode{np.ndarray} (N-dimensional array, или многомерный массив NumPy). Это фундаментальный тип данных в научных вычислениях на Python. В отличие от стандартных списков Python (\mycode{list}), которые могут содержать объекты любых типов и хранятся в памяти в виде разбросанных ссылок, массив \mycode{np.ndarray} устроен принципиально иначе:
\begin{itemize}
    \item \textbf{Однородность (homogeneous):} все элементы массива строго одного типа (обычно вещественные числа двойной точности \mycode{float64}).
    \item \textbf{Непрерывность в памяти (C-contiguous --- расположение элементов в памяти в виде единого последовательного блока, как в языках C и Fortran):} элементы лежат в оперативной памяти единым непрерывным блоком.
    \item \textbf{Векторизация (vectorization --- параллельное выполнение одной математической операции над всеми элементами массива):} математические операции над массивами выполняются на уровне скомпилированного C-кода без медленных циклов Python. Например, выражение \mycode{E - 5.0} мгновенно вычтет число $5.0$ из всех элементов массива параллельно.
\end{itemize}

Именно благодаря свойствам \mycode{np.ndarray} мы можем осуществлять сложнейшую обработку ТГц-сигналов (вычисление средних, наложение окон, преобразование Фурье) в реальном времени со скоростью сотен кадров в секунду!

\subsection{Шаг 1: Удаление постоянного смещения (DC offset)}

Любой реальный прибор вносит в измеряемый сигнал небольшое аппаратное смещение нуля (постоянную составляющую, или DC offset, от англ. direct current offset --- постоянное паразитное напряжение на выходе АЦП (аналого-цифрового преобразователя)) — постоянный вертикальный сдвиг всей ТГц-осциллограммы вверх или вниз относительно истинного нуля. Этот сдвиг возникает из-за тепловых наводок в приемной электронике.

Если применить БПФ (быстрое преобразование Фурье) к сигналу с таким смещением, алгоритм воспримет этот сдвиг как волну нулевой частоты. В результате на спектре возникнет ложный гигантский пик на частоте $0$~ТГц, который скроет от нас полезные данные на низких частотах.

Рецепт исправления предельно прост: нужно вычислить среднее значение сигнала и вычесть его из каждой точки. 

Создай в корне своего проекта новый файл \mycode{spectrum.py} (правой кнопкой мыши по главной папке $\rightarrow$ \mycode{New} $\rightarrow$ \mycode{Python File} $\rightarrow$ назови его \mycode{spectrum}).

В самом начале файла импортируй библиотеку NumPy:
\begin{minted}{python}
import numpy as np
\end{minted}

И добавь туда нашу первую расчетную функцию с явным указанием типов:
\begin{minted}{python}
def remove_dc(E: np.ndarray) -> np.ndarray:
    """
    Удаляет постоянную составляющую (DC offset) из сигнала E,
    вычитая его среднее арифметическое значение.
    """
    return E - np.mean(E)
\end{minted}

\subsubsection{Разбор кода: Функция \mycode{remove_dc}}

Разберём синтаксис и математические операции в нашей первой функции очистки сигнала:
\begin{syntaxdesc}
    \item[Вычисление среднего значения \mycode{np.mean(E)}] Функция из библиотеки NumPy рассчитывает среднее арифметическое всех элементов переданного массива \mycode{E} (то есть суммирует все отсчёты поля и делит их на общее количество точек).
    \item[Векторизованное вычитание \mycode{E - np.mean(E)}] Классический пример векторизации в NumPy. Из массива (вектора) вычитается одно число (скаляр) --- среднее значение. NumPy автоматически выполняет это действие для каждого элемента массива параллельно, без явных циклов.
\end{syntaxdesc}
Применяемые здесь аннотации типов (\mycode{E: np.ndarray} и возвращаемый тип \mycode{-> np.ndarray}) сообщают среде разработки, что на вход ожидается массив NumPy, и на выходе мы также получаем обработанный массив.

\subsection{Шаг 2: Оконная фильтрация (Гауссово окно)}

Терагерцовый импульс представляет собой очень короткий всплеск поля длительностью менее $1$--$2$~пс. Однако запись сигнала длится значительно дольше (например, порядка $100$~пс). В «хвосте» сигнала полезного импульса уже нет — там регистрируются только случайные шумы прибора и эхо-сигналы (паразитные переотражения ТГц-волны внутри деталей прибора).

Чтобы аккуратно вырезать полезный импульс и плавно подавить шум на краях, мы применим оконную фильтрацию — умножим наш временной сигнал на оконную функцию $W(t)$ (Гауссово окно), которая описывается классическим уравнением:
\begin{equation}
W(t) = \exp\left( -\frac{(t - t_{\text{peak}})^2}{2\sigma^2} \right)
\end{equation}
где $t_{\text{peak}}$ — положение временного центра (пика) ТГц-импульса, а $\sigma$ — полуширина окна, определяющая скорость затухания функции к краям записи. 

Ширину этого окна во временной области мы зададим в файле конфигурации \mycode{config.py} параметром \mycode{SIGMA_PS}. Открой свой файл \mycode{config.py} и добавь в его конец следующую строку: \mycode{SIGMA_PS = 20.0} (это стандартное значение полуширины окна $\sigma = 20.0$~пс, которое позволит аккуратно локализовать ТГц-импульс).

Перед генерацией окна необходимо автоматически найти положение главного пика импульса $t_{\text{peak}}$, чтобы точно отцентрировать окно по временной шкале:

\begin{listing}[H]
	\caption{Центрирование и наложение Гауссова окна (\mycode{spectrum.py})}
	\begin{minted}{python}
import numpy as np
import config

def find_peak_center(E: np.ndarray) -> int:
    """Возвращает индекс максимального по модулю значения (центр импульса)."""
    return int(np.argmax(np.abs(E)))

def apply_gaussian_window(t: np.ndarray, E: np.ndarray, sigma_ps: float = None, peak_idx: int = None) -> tuple[np.ndarray, np.ndarray]:
    """
    Создает Гауссово окно с заданной шириной sigma_ps (в пикосекундах) 
    вокруг временного пика и плавно сводит сигнал к нулю по краям.
    """
    if sigma_ps is None:
        sigma_ps = config.SIGMA_PS
        
    if peak_idx is None:
        peak_idx = find_peak_center(E)
        
    t_peak = t[peak_idx]
    
    # Математическая формула Гауссовой кривой: exp(- (t - t_peak)^2 / (2 * sigma^2))
    window = np.exp(-0.5 * ((t - t_peak) / sigma_ps) ** 2)
    
    E_windowed = E * window
    return E_windowed, window
	\end{minted}
\end{listing}



\subsubsection{Разбор кода: Оконная фильтрация}

Давай внимательно разберем заголовок и структуру функции \mycode{apply_gaussian_window}. Здесь мы видим несколько новых для нас конструкций современного Python:
\begin{syntaxdesc}
    \item[Ширина по умолчанию \mycode{sigma_ps: float = None} (строка 7)] Конструкция \mycode{= None} указывает, что параметр \mycode{sigma_ps} является опциональным (необязательным для передачи при вызове функции). Если при вызове функции мы не укажем ширину окна (например, вызовем \mycode{apply_gaussian_window}), Python автоматически присвоит переменной \mycode{sigma_ps} специальное пустое значение \mycode{None}. Внутри функции мы проверяем это условие (\mycode{if sigma_ps is None: sigma_ps = config.SIGMA_PS}) и подставляем стандартное (заданное по умолчанию) значение из нашего конфигурационного файла. Это очень красивый шаблон проектирования (типовое архитектурное решение): функция остается гибкой (мы можем передать любую ширину), но при этом имеет надежное стандартное поведение «из коробки».
    
    \item[Положение пика по умолчанию \mycode{peak_idx: int = None} (строка 7)] Этот параметр позволяет вручную передать индекс максимума ТГц-импульса. Если он равен \mycode{None}, функция автоматически определяет пик по входному сигналу \mycode{E} через функцию \mycode{find_peak_center}. Зачем это нужно? Когда исследуемый образец вносит очень сильное затухание в ТГц-излучение (например, при исследовании сильно поглощающих материалов), полезный сигнал становится настолько слабым и зашумленным, что автоматический алгоритм может «промахнуться» мимо реального центра импульса и зацепиться за случайный шумовой всплеск на краях. В этом случае мы можем точно определить центр импульса по стабильному и яркому опорному сигналу фона (background) и принудительно передать этот индекс через параметр \mycode{peak_idx} для обоих сигналов, гарантируя абсолютно идентичное и ровное наложение оконного фильтра.
    
    \item[Множественный возврат \mycode{-> tuple[np.ndarray, np.ndarray]} (строка 7)] Конструкция \mycode{tuple} (кортеж / упорядоченный неизменяемый набор данных) указывает на то, что функция возвращает не одиночное значение, а кортеж из двух элементов. В квадратных скобках \mycode{[np.ndarray, np.ndarray]} подробно расписано, что и первый, и второй элементы этого кортежа являются массивами NumPy. В конце функции мы видим строчку \mycode{return E_windowed, window}. Это позволяет нам при вызове функции элегантно распаковать результат в две независимые переменные одной строкой: \mycode{E_filtered, W = apply_gaussian_window(t, E)}.
\end{syntaxdesc}

Кроме борьбы с эхо-сигналами, у оконного фильтра есть еще один важный практический эффект — контролируемое сглаживание спектра. В реальном ТГц-эксперименте на спектрах часто видны очень резкие и глубокие провалы — узкие линии поглощения (например, молекулами водяного пара в воздухе). Применение Гауссова окна сглаживает спектральную кривую: узкие линии поглощения и мелкие шумовые осцилляции становятся менее глубокими и размываются. Это помогает убрать лишнюю «изрезанность» графика и получить более чистую, плавную огибающую спектра пропускания нашей решетки.

\warning{Вычитание постоянной составляющей (DC offset) обязательно должно выполняться \textbf{до} наложения окна! Если сначала умножить сигнал со смещением на окно, постоянный сдвиг превратится в плавный Гауссов «горб» во временной области. Такой искусственный сигнал уже невозможно отфильтровать простым вычитанием константы.}

\subsection{Шаг 3: Расчет быстрого преобразования Фурье (БПФ)}

Теперь мы готовы перевести временной ТГц-импульс в спектр. В Python Быстрое преобразование Фурье реализовано в стандартной библиотеке \mycode{scipy.fft}. Так как наше поле $E(t)$ состоит из чисто вещественных (действительных) чисел, мы используем специализированную функцию \mycode{rfft} (Real FFT --- преобразование Фурье для вещественного сигнала) вместо стандартной комплексной \mycode{fft} (Fast Fourier Transform --- быстрое преобразование Фурье для комплексного сигнала). Это экономит вычислительные ресурсы и автоматически убирает избыточные отрицательные частоты.

\note{\textbf{Важная деталь о названии «Real FFT»:} у начинающих программистов часто возникает путаница. Префикс \mycode{r} (от Real — вещественный) в названии функции \mycode{rfft} относится исключительно к \textbf{входным данным} (то есть на вход подается массив вещественных чисел \mycode{E}, у которых мнимая часть равна нулю). Однако сам \textbf{результат работы функции} (выходной спектр) остается \textbf{комплексным}! Каждая спектральная гармоника по-прежнему имеет как амплитуду, так и фазу, поэтому возвращаемый массив \mycode{spectrum_complex} состоит из комплексных чисел (\mycode{complex128}). Чтобы избавиться от фазовой составляющей и получить привычный вещественный амплитудный спектр, мы обязаны взять модуль комплексных чисел с помощью функции \mycode{np.abs(spectrum_complex)}.}

Давай строго зафиксируем физические и математические обозначения, чтобы избежать путаницы между временной и частотной областями:
\begin{itemize}
    \item $E(t)$ — мгновенная напряженность электрического поля ТГц-импульса как функция времени (наш исходный вещественный сигнал, в коде — массив \mycode{E}).
    \item $E(\nu)$ — комплексный спектр поля как функция частоты $\nu$ (в коде — массив \mycode{spectrum_complex}). Каждая спектральная компонента имеет вид $E(\nu) = A(\nu) \cdot e^{i \varphi(\nu)}$, где $\varphi(\nu)$ — спектральная фаза.
    \item $A(\nu) = |E(\nu)|$ — вещественный амплитудный спектр (в коде — массив \mycode{amplitude_spectrum}), который получается как абсолютное значение (модуль) комплексного спектра $E(\nu)$. Именно эту величину мы строим на спектральных графиках.
\end{itemize}

Частотную шкалу мы получим с помощью функции \mycode{rfftfreq}. Важно помнить, что ось частот, генерируемая этой функцией, содержит ровно $N/2 + 1$ точек (где $N$ --- длина исходного временного сигнала). Так как физический сигнал вещественный, его отрицательные частоты не несут новой информации (спектр зеркально симметричен). Функция \mycode{rfft} отбрасывает эту избыточную отрицательную половину, экономя оперативную память и ускоряя расчеты в два раза!

\begin{listing}[H]
	\caption{Вычисление спектра амплитуд (\mycode{spectrum.py})}\label{lst:compute_spectrum}
	\begin{minted}{python}
import numpy as np
from scipy.fft import rfft, rfftfreq

def compute_spectrum(t: np.ndarray, E: np.ndarray) -> tuple[np.ndarray, np.ndarray]:
    """
    Вычисляет спектр амплитуд ТГц-импульса.
    Возвращает ось частот (в ТГц) и вещественный амплитудный спектр.
    """
    N = len(E)
    dt = t[1] - t[0]  # Шаг временной сетки в пикосекундах
    
    # Вычисляем преобразование Фурье для вещественного сигнала
    spectrum_complex = rfft(E)
    
    # Модуль комплексного числа дает амплитуду спектральной гармоники
    amplitude_spectrum = np.abs(spectrum_complex)
    
    # Генерация шкалы частот.
    # Так как dt задан в пикосекундах, частоты freqs будут в ТГц
    freqs = rfftfreq(N, d=dt)
    
    return freqs, amplitude_spectrum
	\end{minted}
\end{listing}



\subsection{Шаг 4: Усреднение сигналов в спектральной области}

При проведении ТГц-эксперимента для каждого угла поляризатора обычно записывают несколько повторных измерений (проходов) сигнала, чтобы повысить точность.

\warning{Ни в коем случае не усредняй исходные временные сигналы нескольких проходов перед расчетом БПФ! \\
Из-за незначительного случайного временного дрожания прибора (фазового джиттера (временной нестабильности срабатывания триггера оцифровки сигналов)) ТГц-импульсы в разных проходах могут быть микроскопически сдвинуты друг относительно друга по оси времени. Если их усреднить во временной области, они частично погасят друг друга. На итоговом спектре это приведет к ложному падению амплитуд на высоких частотах.\\
\textbf{Правильный рецепт:} сначала рассчитай БПФ-спектр (спектр быстрого преобразования Фурье) для каждого прохода по отдельности, возьми модуль (амплитуду), и только после этого усредняй полученные амплитудные графики на частотной сетке.}

\subsection{Шаг 5: Вычисление пропускания по мощности}

Коэффициент пропускания $T(\nu)$ исследуемой решетки по мощности (power transmission) в ТГц-диапазоне показывает, какая часть мощности излучения прошла сквозь нее на каждой частоте. 

В отличие от традиционной оптической спектроскопии, где детекторы реагируют только на интенсивность света (усредненную по времени мощность), спектроскопия терагерцового разрешения во времени (TDS) обладает уникальной особенностью: наши датчики регистрируют непосредственно \textbf{амплитуду электрического поля} $E(t)$ волны! 

По законам электродинамики, мгновенная мощность электромагнитной волны пропорциональна квадрату ее амплитуды ($I \propto |E|^2$). Соответственно, чтобы рассчитать пропускание по мощности $T(\nu)$ (отношение прошедшей мощности к падающей), мы должны разделить интенсивность волны, прошедшей через образец, на интенсивность волны в воздухе (фоне), что математически выражается как квадрат отношения их спектральных амплитуд:
\begin{equation}
T(\nu) = \frac{I_{\text{sample}}(\nu)}{I_{\text{bg}}(\nu)} = \left| \frac{E_{\text{sample}}(\nu)}{E_{\text{bg}}(\nu)} \right|^2 = \left( \frac{A_{\text{sample}}(\nu)}{A_{\text{bg}}(\nu)} \right)^2
\end{equation}

Напишем функцию \mycode{calculate_transmission}, которая объединяет все шаги в единый конвейер (последовательную цепочку обработки):

\begin{listing}[H]
	\caption{Конвейер расчета спектра пропускания (\mycode{spectrum.py})}\label{lst:calculate_transmission}
	\begin{minted}{python}
import numpy as np

def calculate_transmission(
        t_sig: np.ndarray, E_sig: np.ndarray,
        t_bg: np.ndarray, E_bg: np.ndarray,
) -> tuple[np.ndarray, np.ndarray, np.ndarray, np.ndarray]:
    """Полный цикл обработки: от сырого ТГц-сигнала до спектра пропускания по мощности."""
    
    # 1. Удаление постоянного смещения
    E_sig_clean = remove_dc(E_sig)
    E_bg_clean = remove_dc(E_bg)
    
    # Поиск центра пика по опорному сигналу (фона) как золотой стандарт
    peak_idx = find_peak_center(E_bg_clean)
    
    # 2. Наложение временного Гауссова окна (центрированного по пику опорного сигнала)
    E_sig_win, _ = apply_gaussian_window(t_sig, E_sig_clean, peak_idx=peak_idx)
    E_bg_win, _ = apply_gaussian_window(t_bg, E_bg_clean, peak_idx=peak_idx)
    
    # 3. Расчет БПФ
    freqs, spec_sig = compute_spectrum(t_sig, E_sig_win)
    _, spec_bg = compute_spectrum(t_bg, E_bg_win)
    
    # 4. Расчет пропускания по мощности (T = |E_sig / E_bg|^2)
    # Используем np.maximum для предотвращения деления на ноль на краях диапазона
    spec_bg_safe = np.maximum(spec_bg, 1e-10)
    transmission = (spec_sig / spec_bg_safe) ** 2
    
    return freqs, spec_sig, spec_bg, transmission
	\end{minted}
\end{listing}

\subsubsection{Разбор кода: Спектральные преобразования}

Давай рассмотрим основные элементы конвейера:
\begin{syntaxdesc}
    \item[Вычисление БПФ \mycode{rfft(E)} (строка 12 в Листинге~\ref{lst:compute_spectrum})] Стандартная функция из \mycode{scipy.fft}, которая преобразует массив вещественных чисел \mycode{E} во временной области в массив комплексных спектральных гармоник.
    \item[Извлечение амплитуды \mycode{np.abs(spectrum_complex)} (строка 15 в Листинге~\ref{lst:compute_spectrum})] Отбрасывает фазовую информацию, оставляя только физическую амплитуду гармоники.
    \item[Центрирование по опорному сигналу \mycode{peak_idx = find_peak_center(E_bg_clean)} (строка 10 в Листинге~\ref{lst:calculate_transmission})] Обрати внимание на очень важную физическую особенность: мы центрируем временное окно по пику опорного сигнала (фона) для обоих измерений. Наш образец представляет собой тончайший проволочный поляризатор без подложки. Он настолько тонкий (толщина проволоки всего около $11$~мкм), что физически не вносит никакой задержки во времени (фазового сдвига) в прошедшую волну относительно воздуха. Если бы мы измеряли толстый образец-диэлектрик (например, пластину кремния), то из-за преломления ($n > 1$) волна в нем распространялась бы медленнее ($v = c/n$), и прошедший ТГц-импульс пришел бы с существенной временной задержкой. В этом случае положение окна для образца пришлось бы сдвигать. Но для нашей тонкой проволочной решетки такое смещение отсутствует, что позволяет использовать индекс пика фона \mycode{peak_idx} как универсальный золотой стандарт центрирования для обоих файлов.
    \item[Безопасное деление \mycode{np.maximum(spec_bg, 1e-10)} (строка 22 в Листинге~\ref{lst:calculate_transmission})] Функция \mycode{np.maximum} поэлементно сравнивает спектр фона с исчезающе малым числом \mycode{1e-10} (запись в научном формате Python, эквивалентная числу $10^{-10}$ или $0.0000000001$) и заменяет любые околонулевые значения шума на это число. Это гарантирует, что мы никогда не столкнемся с критической ошибкой деления на ноль, даже если спектральная амплитуда фона на краях чувствительности прибора упадет до абсолютного нуля.
\end{syntaxdesc}

\begin{minitaskbox}
	\textbf{Мини-задание: Численный эксперимент с шириной Гауссова окна} \\
	Для того чтобы наглядно прочувствовать физическую связь между временным окном и его спектральным представлением, тебе предлагается провести простой численный эксперимент. Напиши небольшой тестовый скрипт, в котором:
	\begin{enumerate}
		\item Загрузи тестовый ТГц-сигнал с помощью разработанной ранее функции \mycode{load_tds_data}.
		\item Примени к сигналу оконную фильтрацию \mycode{apply_gaussian_window} последовательно с тремя разными значениями полуширины окна $\sigma_t$: $5.0$~пс (очень узкое окно), $20.0$~пс (стандартное сбалансированное окно) и $60.0$~пс (очень широкое окно, практически не затухающее к краям).
		\item Рассчитай амплитудный спектр для каждого из этих трёх случаев с помощью \mycode{compute_spectrum}.
		\item Построй графики исходного и заокненных сигналов во временной области, а также графики их нормированных амплитудных спектров на одном рисунке с помощью библиотеки \mycode{matplotlib}.
	\end{enumerate}
	Сравни результаты. Посмотри, как очень узкое временное окно ($5.0$~пс) сглаживает все спектральные особенности, превращая спектр в одну плавную широкую куполообразную линию, в то время как широкое окно ($60.0$~пс) сохраняет глубокие и резкие линии поглощения атмосферным водяным паром, но оставляет в спектре много мелкого «шума». Сделай вывод о том, как правильно искать баланс при фильтрации ТГц-сигналов.
\end{minitaskbox}

\subsection{Пакетная обработка спектров (Batch Processing)}\label{ssec:process_dataset}

Теперь, когда все сырые данные разложены по полочкам в нашей СУБД, напишем функцию для их автоматической пакетной обработки. Функция должна пройтись по всем уникальным углам, рассчитать спектры пропускания для каждой реплики (строго попарно!), а затем корректно усреднить их по частоте.

Добавим функцию \mycode{process_dataset} в спектральный модуль \mycode{spectrum.py} (полный листинг модуля находится в \textbf{\textit{Приложении~\ref{app:spectrum}}}):

\refstepcounter{listing}\label{lst:process_dataset}
\noindent\textbf{Листинг~\thelisting~---~Пакетная обработка базы данных (\mycode{spectrum.py})}
\begin{minted}{python}
def process_dataset(data_store: dict) -> dict:
    # Находим все уникальные пары углов (angle1, angle2)
    angle_pairs = set()
    for key in data_store.keys():
        if key[0] == 'signal_raw':
            angle_pairs.add((key[1], key[2]))

    # Для каждой пары углов проводим вычисления
    for angle1, angle2 in sorted(angle_pairs):
        transmissions = []
        bg_spectra = []
        freqs_common = None

        reps = [
            key[3] for key in data_store.keys()
            if key[0] == 'signal_raw'
            and key[1] == angle1
            and key[2] == angle2
        ]

        for rep in sorted(reps):
            sig_key = ('signal_raw', angle1, angle2, rep)
            bg_key = ('bg_raw', angle1, angle2, rep)

            if bg_key in data_store:
                t_sig, E_sig = data_store[sig_key]
                t_bg, E_bg = data_store[bg_key]

                # Рассчитываем пропускание конкретной реплики
                freqs, spec_sig, spec_bg, transmission = calculate_transmission(
                    t_sig, E_sig, t_bg, E_bg
                )

                transmissions.append(transmission)
                bg_spectra.append(spec_bg)
                if freqs_common is None:
                    freqs_common = freqs

\end{minted}
\newpage
\begin{minted}[firstnumber=39, autogobble=false]{python}
        if transmissions:
            # Математическое усреднение для компенсации дрейфа
            trans_avg = np.mean(transmissions, axis=0)
            bg_avg = np.mean(bg_spectra, axis=0)

            # Сохраняем агрегированные данные обратно в СУБД
            data_store[('spec_avg', angle1, angle2, None)] = (
                freqs_common, trans_avg
            )
            data_store[('spec_bg_avg', angle1, angle2, None)] = (
                freqs_common, bg_avg
            )

    return data_store
\end{minted}

\subsubsection{Разбор кода: Пакетный обработчик}

\begin{syntaxdesc}
    \item[Сбор пар углов \mycode{angle_pairs = set()} (строка 3)] Множество \mycode{set()} автоматически отбрасывает дубликаты. Поэтому даже если в базе есть 10 реплик для угла \mycode{(10.0, 0.0)}, в \mycode{angle_pairs} он попадёт ровно один раз.
    \item[Фильтрация реплик \mycode{reps = [key[3] for key in ...]} (строки 14--19)] Генератор списков (list comprehension --- компактный способ создания списков) отбирает из базы данных номера всех реплик для текущей пары углов. Номер реплики --- это четвёртый элемент ключа \mycode{key[3]}.
    \item[Проверка наличия фона \mycode{if bg_key in data_store} (строка 25)] Гарантирует, что мы не вызовем функцию \mycode{calculate_transmission} с отсутствующими фоновыми данными. Если фоновые файлы отсутствуют или не загрузились корректно, реплика просто пропускается.
    \item[Арифметическое усреднение \mycode{np.mean(transmissions, axis=0)} (строка 41)] Функция \mycode{np.mean} с параметром \mycode{axis=0} усредняет набор одномерных массивов-спектров по первому измерению (<<по колонке>>). То есть для каждого частотного бина (частотной ячейки / отсчета) вычисляется среднее значение по всем репликам одновременно.
    \item[Тег усредненных данных \mycode{('spec_avg', angle1, angle2, None)} (строка 45)] Четвёртый элемент \mycode{None} вместо конкретного номера реплики явно сигнализирует: эта запись представляет собой итоговый усредненный спектр для данной пары углов.
\end{syntaxdesc}

\note{Благодаря векторным операциям NumPy, усреднение семейства спектров выполняется мгновенно с помощью \mycode{np.mean(..., axis=0)}. Это позволяет обрабатывать огромные экспериментальные массивы за доли секунды.}

\subsection{Запуск и проверка расчетного модуля}\label{ssec:spectrum_run}

Давай проверим наш новый модуль \mycode{spectrum.py} в деле! В конце полного файла (его полный листинг приведен в \textbf{\textit{Приложении~\ref{app:spectrum}}}) мы написали встроенный блок самотестирования. Он делает следующее:
\begin{enumerate}
    \item Загружает реальные файлы сигналов образца и фона из твоего лабораторного архива с помощью ранее разработанного модуля \mycode{data_loader.py}.
    \item Очищает их от постоянного смещения с помощью функции \mycode{remove_dc} и аккуратно накладывает Гауссово окно через \mycode{apply_gaussian_window}.
    \item Вычисляет спектры и пропускание по мощности с помощью единого сквозного конвейера \mycode{calculate_transmission}.
    \item Строит красивый двухпанельный проверочный график и сохраняет его в корневой папке проекта в файл \mycode{test_spectrum_dsp.png}.
\end{enumerate}

Выполни запуск нашего расчетного модуля:
\note{Кликни правой кнопкой мыши в любом месте внутри редактора кода файла \mycode{spectrum.py} и выбери пункт \textbf{Run 'spectrum'}.}

В окне вывода PyCharm (Зона 4) ты должен увидеть следующее подтверждение:
\begin{minted}{text}
Запуск встроенного блока самотестирования модуля spectrum.py на реальных данных...
  - Среднее значение до удаления DC: 4.512330761623351e-05
  - Среднее значение после удаления DC: -1.341103761109312e-20
  - Координата найденного центра пика: 51.25 пс
Тестирование успешно завершено! График сохранен в файл: test_spectrum_dsp.png
\end{minted}

\warning{Обрати внимание на среднее значение после удаления постоянной составляющей: оно равно числу порядка $10^{-20}$ в научных обозначениях Python (буква \mycode{e} обозначает «десять в степени», то есть $-1.34 \times 10^{-20}$). Это фактически абсолютный физический ноль на языке компьютера! Наш код сработал безукоризненно.}

Открой файл \mycode{test_spectrum_dsp.png}, который появился в корневой папке твоего проекта (в проводнике PyCharm слева). На нем ты увидишь две панели:
\begin{itemize}
    \item \textbf{На верхней панели} отображаются исходный и очищенный ТГц-сигнал, плавная колоколообразная кривая Гауссова окна и итоговый заокненный сигнал. Посмотри, как плавно затухает сигнал на краях, уводя в ноль случайные шумы!
    \item \textbf{На нижней панели} показаны амплитудные спектры фона (воздуха) и образца, а также итоговый коэффициент пропускания по мощности $T(\nu)$. Обрати внимание на резкие узкие провалы в спектрах вблизи $1.15$~ТГц, $1.41$~ТГц и $1.67$~ТГц --- это физические линии поглощения ТГц-излучения водяным паром, содержащимся в воздухе твоей комнаты!
\end{itemize}

\subsection{Полный листинг модуля}

Полный, готовый к работе исходный код расчетного модуля \mycode{spectrum.py} (с решенными проверочными блоками и встроенным блоком самотестирования на реальных данных) вынесен в \textbf{\textit{Приложение~\ref{app:spectrum}}}. Ты всегда можешь использовать его для сверки со своим кодом.

\subsection{Итоги раздела}
В этом разделе мы научились применять базовые цифровые методы для обработки экспериментальных ТГц-сигналов:
\begin{enumerate}
	\item Очистили сигнал от постоянного смещения с помощью функции \mycode{remove_dc}.
	\item Научились с помощью функции \mycode{apply_gaussian_window} для Гауссова окна плавно вырезать полезную часть импульса.
	\item Использовали функцию \mycode{compute_spectrum} для перехода в частотную область.
	\item Объединили все операции в конвейер (последовательную цепочку обработки) \mycode{calculate_transmission} для расчёта спектра пропускания.
	\item Написали пакетный обработчик \mycode{process_dataset}, который автоматически проходит по всем углам, рассчитывает спектры попарно для каждой реплики и усредняет результаты.
\end{enumerate}

Теперь вся база данных полностью обработана и готова к визуализации. В следующем разделе мы построим интерактивный графический интерфейс для поштучного просмотра семейства спектров и построения всей угловой зависимости пропускания!
